% The appendices of the notes: Bilaga A is the method, written for the
% students; then one chapter per backed claim.  PROJECT: the headings
% below are Swedish; translate them for a deck in another language.

\chapter{Att kontrollera ett påstående}
\label{app:verifiera}

% Written for the student, not for the teacher.  Open by listing the
% deck's claims that cannot be checked at the keyboard, each with a \cref
% to the section that makes it, then say which appendix takes which.
\section{Gör påståendet till en fråga}
% A claim can be checked only once it is clear what would show it wrong.
\section{Välj rätt sorts källa}
% Primary source for origin, normative documentation for what a tool
% does, reviews for whether something is established, single studies for
% one result.
\section{Sök så att du kan ha fel}
% Subject-driven, not author-driven; the same queries in every database;
% counter-queries; known documents when indexing fails.
\section{Läs i fulltext, och skriv ner hur}
% A hit is not evidence: quote with a page number, or say the source was
% read in abstract only.  Then the provenance fields CLAIM, FOUND-VIA,
% PICKED, QUOTE, VERIFIED, COUNTER, DATE.
\section{Dra slutsatsen --- med förbehåll}
% What the search supports, at the strength it supports it.
\section{En anmärkning om språkmodeller}

% ---------------------------------------------------------------- claim
% One chapter per backed claim.  The title is the question, so the reader
% knows what is being answered before reading a word of it.
\chapter{Question the search answers?}
\label{app:claim}
\chapterprecis{Författaren har ännu inte granskat resultaten i den här
  bilagan i sin helhet.}

% Where the claim is made in the deck, what exactly is claimed (often two
% claims: establishment and origin), and the question that follows.  Then
% a road map: Metod, Resultat, Diskussion, Slutsats, Fortsatt arbete, and
% the hit list last.

\section{Metod}
\label{sec:claim-metod}

% Date, which databases answered and which refused, what the supporting
% queries look for and what the counter-queries look for, then the query
% table.  Known documents fetched outside the databases are named here
% with the reason.

\begin{table}[!htb]
  \caption{Sökfrågor och antal träffar per databas, DATE, verktyget
    \texttt{scholar}.  S = stödfråga, M = motfråga.  Databaser: OpenAlex
    (OA; inte \foreignlanguage{english}{open access}), DBLP, IEEE Xplore
    (IEEE), Web of Science (WoS), Scopus.  Högst 40 träffar hämtades per
    databas och fråga, så 40 betyder \enquote{minst 40}.}%
  \label{tab:sok-queries-claim}
  \centering\footnotesize
  \begin{tabular}{@{}lp{0.42\linewidth}rrrrr@{}}
    \toprule
       & Nyckelord & OA & DBLP & IEEE & WoS & Scopus \\
    \midrule
    S1 & the supporting query, in the boolean form sent to the databases
       & 40 & 0 & 13 & 16 & 26 \\
    \midrule
    M1 & the counter-query, with the same concepts plus critique,
         limitations, \enquote{considered harmful}
       & 12 & 0 & 2 & 6 & 8 \\
    \bottomrule
  \end{tabular}
\end{table}

% How many unique records, how they were screened, how many were read in
% full text and how many only in abstract.

\section{Resultat}
\label{sec:claim-resultat}
% What the sources say, supporting and qualifying alike, with quotes.

\section{Diskussion och begränsningar}
\label{sec:claim-diskussion}

\section{Slutsats}
\label{sec:claim-slutsats}
% The first sentence answers the chapter's question.

\section{Fortsatt arbete}
\label{sec:claim-fortsatt}
% What this search left undone, and what the field would have to study
% for a surer answer.

\section{Träffar som rör påståendet}
\label{sec:claim-traffar}
% The hit list is generated by the search tool and \input here.
\input{litteratursokning/claim-full}
